\section{Multimodal Guardrails}
\label{sec:multimodal}

Multimodal large language models—systems that accept inputs in more than one
modality (typically text and image, with audio and video emerging)—introduce
qualitatively new safety challenges.
This section reviews the attack surface, the proposed defences, and the
confirmed and unconfirmed findings from the literature.

\subsection{The Multimodal Attack Surface}

\subsubsection{Cross-Modal Jailbreaks}

Cross-modal jailbreaks exploit the gap between where safety alignment is
applied (predominantly the text backbone) and where harmful content can enter
(any input modality).
Several attack classes have been documented:

\paragraph{Typographic attacks}
Embedding text into images allows harmful instructions to enter the model's context
through the vision encoder rather than the text tokeniser, potentially bypassing
text-focused filters~\cite{gong2023figstep}.
The FigStep attack converts harmful text into images of step-by-step instructions,
exploiting VLMs' tendency to follow visually presented instructions.

\paragraph{Adversarial image perturbations}
Imperceptible perturbations applied to images can steer VLM outputs toward
attacker-chosen strings~\cite{liu2023queryrelevant}.
These attacks target the visual encoder (e.g.\ CLIP) rather than the language
backbone.

\paragraph{Indirect prompt injection via media}
An attacker can embed adversarial instructions in an image or document that the user
shares with the model, hijacking the model's behaviour without the user's knowledge.
AgentPoison~\cite{chen2024agentpoison} extends this vector to agentic LLM systems,
showing that poisoning a model's external memory or knowledge base can redirect
agent behaviour without any modification to the model's weights.
A clarification on scope is warranted: AgentPoison is evaluated primarily on
text-only agentic systems with multimodal retrieval (i.e.\ systems that retrieve
and process image or document content but whose core reasoning backbone is a
text LLM), rather than on vision-language models where the image encoder is
part of the model's inference graph.
The attack surface distinction matters: in a text-only agent with multimodal
retrieval, the attack operates through the retrieved context that is transcribed
or described for the text LLM, whereas in a true VLM the adversarial content
can be embedded directly in the image encoder's input.
AgentPoison is included under multimodal attacks because it demonstrates the
retrieval-augmentation attack vector that applies when VLMs are deployed in
agentic contexts with external memory or knowledge bases.

\paragraph{Text-channel attacks on multimodal models}
A counterintuitive finding reported across multiple papers~\cite{zhao2024evaluating}
is that text-only jailbreaks applied to MLLMs---without any image input---achieve
high attack success rates, suggesting that multimodal fine-tuning may weaken the
text-backbone's safety alignment.
These specific numerical claims (average ASR of 50.5\% across open-source models)
did not survive adversarial verification in this review.
The \emph{qualitative direction} of this finding---that MLLM text channels are
more vulnerable than equivalent text-only LLM baselines---is treated in this
review as an \emph{unverified but plausible hypothesis}, not as an established
fact: it is consistent with the theoretical framing (multimodal fine-tuning
modifies the text backbone's activation distribution) but no study has provided
a reproducible controlled comparison isolating the contribution of multimodal
fine-tuning from confounds such as different training data, safety fine-tuning
budget, or RLHF stage.
This distinction matters for open problem MM1: the uncertainty is about the
\emph{mechanism} (is it fine-tuning per se, or data differences?) as much as the
\emph{magnitude}.

\subsubsection{The Multi-Turn Conversational Attack}

Of the 26 quantitative and architectural claims extracted from the multimodal
guardrails corpus (see Table~\ref{tab:claim_counts}), two survived structured
self-review: the multi-turn conversational priming finding and the cross-modal
attack surface taxonomy.
The only multimodal-relevant attack mechanism that survived all three temporally
separated self-review verification passes in this review concerns multi-turn
conversational context~\cite{li2024llm,russinovich2024greatescalation}.
Multi-turn attacks prime the model through a sequence of benign exchanges that
gradually shift the conversational context before introducing a harmful request.
The Crescendo attack~\cite{russinovich2024greatescalation} reports that this
approach outperforms single-turn jailbreaks by 29--61\% on GPT-4 and 49--71\% on
Gemini-Pro across seven evaluated models.
This claim passed all three verification passes for three specific reasons that
distinguish it from analogous numerical claims that did not survive (e.g.\
WildGuard's ASR reduction from 79.8\% to 2.4\%): (i) Russinovich et al.\ specify
the exact prompting procedure, the seven target models, the evaluated harm
categories, and the ASR measurement methodology in sufficient detail for a reader
to reconstruct the evaluation; (ii) the directional finding (multi-turn priming
outperforms single-turn) is corroborated by an independent replication framing
in~\cite{li2024llm}, which reports consistent multi-turn amplification using a
different attack methodology; and (iii) the percentage-improvement figures are
reported as ranges across models rather than single point estimates, which is
internally consistent with the observed variance across seven models and reduces
the risk of cherry-picking.
By contrast, single-point ASR claims such as WildGuard's 2.4\% failed because the
evaluation baseline, judge model version, and token budget were not fully specified,
making independent reproduction infeasible.

The implication for multimodal guardrails is direct: a guardrail that inspects only
the current turn---whether text or image---is structurally insufficient.
Session-level conversational history must be a first-class input to any safety
classifier intended to be robust against multi-turn attacks.
This architectural requirement has not been incorporated into any of the verified
real-time guardrail systems reviewed in Section~\ref{sec:realtime}.

\subsection{VLM Safety Alignment}

Table~\ref{tab:vlm_defence} compares the three principal VLM defence architectures
reviewed in this section.
All performance figures are as reported by the respective authors and are marked
unverified (U) in accordance with the adversarial verification protocol
(Section~\ref{subsec:verification_protocol}).

\begin{table}[t]
  \caption{Comparison of proposed VLM defence architectures. Performance figures
           are author-reported and \textbf{unverified (U)}. VLM = requires retraining
           the base vision-language model; N/A = not reported.}
  \label{tab:vlm_defence}
  \centering
  \begin{tabular}{p{2.5cm}p{3.0cm}p{2.5cm}p{2.5cm}p{1.5cm}}
    \toprule
    \textbf{System} & \textbf{Design principle} & \textbf{Claimed gain} & \textbf{Retrains VLM?} & \textbf{Status} \\
    \midrule
    SafeVLM~\cite{oh2024safevlm}
      & Plug-in safety projector + tokens + head
      & Safety score $\uparrow$; ASR $\downarrow$ (unspecified)
      & No
      & U \\
    DRESS$^\dagger$~\cite{chen2024dresssafety}
      & NL-feedback RL alignment (instruction-following primary)
      & $+21\%$ harmlessness (benchmark-specific)
      & Yes
      & U \\
    MLLM-Protector~\cite{pi2024mllm}
      & Detect-then-detoxify two-stage adapter
      & ASR $\downarrow$ across multiple attack types
      & No
      & U \\
    \bottomrule
  \end{tabular}
  \par\smallskip
  \noindent\small $^\dagger$ Peripheral inclusion: DRESS is primarily an
  instruction-following alignment paper; it is included because a subset of its
  evaluation measures harmlessness, to contrast training-time approaches with the
  inference-time guardrails above.
  It should not be treated as evidence of a dedicated VLM safety guardrail.
\end{table}

The key cross-system observation is that the two non-retraining approaches
(SafeVLM, MLLM-Protector) are architecturally compatible with deployment pipelines
where the base VLM cannot be retrained, but both lack independently reproducible
evaluations.
DRESS requires full retraining and is included for completeness as a training-time
baseline; it is not a drop-in inference-time guardrail.
None of the three systems incorporates session-level context, limiting their
effectiveness against multi-turn attacks (open problem MM2).

\subsubsection{Does Text Alignment Transfer to Vision?}

A recurring concern across the multimodal safety literature is whether safety
alignment in the text backbone of a VLM transfers to the visual modality.
Several papers report that vision-language fine-tuning weakens safety alignment,
citing specific attack success rate differentials between base LLMs and their
multimodal derivatives.
These quantitative claims did not survive adversarial verification in this review.
The \emph{qualitative} concern---that multimodal training introduces new vulnerability
pathways that text-only safety alignment does not address---is consistent with the
theoretical framing and warrants systematic empirical study with reproducible protocols.

\subsubsection{Proposed Architectures}

\paragraph{SafeVLM}
SafeVLM~\cite{oh2024safevlm} adds three modules to an existing VLM: a safety
projector, safety tokens injected into the input sequence, and a safety
classification head.
The claimed improvements in safety score and reductions in attack success rate did
not survive adversarial verification.
The modular, plug-in architecture---avoiding retraining the full VLM---represents
a practically important design principle.

\paragraph{DRESS}
DRESS~\cite{chen2024dresssafety} is a training-time, model-intrinsic alignment
paper included solely as a contrast baseline to the inference-time guardrails above.
Its peripheral status and the reasons for inclusion are documented in full in the
footnote to Table~\ref{tab:vlm_defence}; no further analysis is offered here to
avoid redundancy.

\paragraph{MLLM-Protector}
MLLM-Protector~\cite{pi2024mllm} adds a lightweight harm-detection adapter and a
response-detoxification module to existing MLLMs.
The two-stage architecture---detect then detoxify---is designed to degrade gracefully:
if the detoxification module introduces hallucinations on safe content, the harm
detector can suppress it selectively.

\paragraph{VLGuard}
VLGuard~\cite{li2024vlguard} addresses VLM safety under low-resource constraints,
proposing a safety fine-tuning approach that achieves alignment improvements with a
small curated dataset rather than large-scale RLHF.
The reported gains in harmlessness on multimodal benchmarks did not survive
author-verified scrutiny in this review---the evaluation dataset and baseline
model version are insufficiently specified for independent reproduction---but the
low-resource framing identifies a practically important deployment scenario:
organisations that cannot afford full-scale safety fine-tuning.

\paragraph{AdaShield}
AdaShield~\cite{wang2024adashield} proposes adaptive shield prompting to defend
against structure-based attacks on multimodal LLMs.
Rather than modifying the model, it prepends dynamically selected safety instructions
that anticipate the attack structure.
This inference-time, model-agnostic approach is architecturally compatible with
deployment pipelines where VLM retraining is infeasible.
Claimed attack success rate reductions did not survive author-verified scrutiny, but
the adaptive prompt construction mechanism is a meaningful design contribution.

\subsubsection{Multimodal Safety Benchmarks}

Systematic evaluation of VLM safety requires standardised benchmarks.
Two principal benchmarks appear in the corpus:
MM-SafetyBench~\cite{liu2024mmsafetybench} evaluates multimodal LLMs across
thirteen safety scenarios using typographic and visual adversarial inputs, revealing
that open-source VLMs are substantially more vulnerable to image-based jailbreaks
than their text-only counterparts; JailBreakV-28K~\cite{luo2024jailbreakv} provides
a large-scale benchmark of 28,000 multimodal jailbreak attempts covering both
image-based and text-based attack vectors, enabling systematic robustness comparison
across models.
Both benchmarks confirm the directional finding that visual inputs represent an
under-defended attack surface, consistent with the theoretical framing in this
section, though the specific quantitative results from individual model evaluations
did not survive author-verified scrutiny due to model-version sensitivity.

\subsection{Audio and Video Safety}
\label{subsec:audiovideo}

Audio and video safety for LLMs is nearly absent from the peer-reviewed literature
as of 2025.
No verified claims about audio-language model safety mechanisms or video-language
model safety mechanisms were found in the three research passes conducted for this
review.
The absence is notable given that audio-capable models (GPT-4o audio mode, Gemini
Live~\cite{google2024gemini}) are in active deployment.

Non-peer-reviewed technical sources were also surveyed.
Anthropic's Claude 3 model card~\cite{anthropic2024claude3} and the Gemini technical
report~\cite{google2024gemini} describe internal safety evaluation processes for
multimodal inputs at a high level but provide no publicly reproducible experimental
protocols or results for audio or video modalities specifically.
These documents were therefore excluded from the primary corpus under the inclusion
criterion requiring empirical or formal evaluation content, but are noted here to
confirm that the gap in academic literature is not filled by available industry
reporting.

Rather than simply documenting this absence, we frame the audio/video safety
problem as a structured research agenda with specific empirical questions that
future work should address in priority order.

\paragraph{Attack-surface characterisation (foundational, should precede defence work)}
The image-safety literature developed attack taxonomies (typographic, adversarial
perturbation, indirect injection, text-channel) before converging on defence
architectures.
An analogous characterisation is needed for audio and video.
Specific open questions:
\begin{itemize}
  \item \textit{Adversarial audio perturbations}: Can imperceptible perturbations
        to speech inputs cause audio-language models to produce harmful outputs,
        analogously to adversarial image attacks on VLMs?
        No empirical study has evaluated this question for any production
        audio-language model.
  \item \textit{Steganographic speech}: Can harmful instructions be encoded in
        audio signals at volumes or frequencies inaudible to humans but
        parseable by speech encoders?
        The image analogue (typographic embedding) is the most effective attack
        class in the image literature; the audio equivalent has not been studied.
  \item \textit{Multi-turn audio priming}: The confirmed finding that multi-turn
        conversational priming amplifies jailbreaks in text-only models
        (Section~\ref{subsec:streaming_gap}) predicts that audio-turn priming
        may have the same effect.
        Whether the dynamics of audio conversation (prosody, speaker identity,
        temporal rhythm) accelerate or retard priming is an open empirical
        question.
  \item \textit{Video-specific attack vectors}: Video inputs add temporal
        structure absent from both image and audio.
        Cross-frame adversarial perturbations, visual scene transitions, and
        embedded text in video frames all represent attack surfaces without
        established empirical characterisation.
\end{itemize}

\paragraph{Defence transfer evaluation (second priority)}
Once attack surfaces are characterised, the critical empirical question is whether
existing image-safety defences transfer to audio and video.
Specific testable hypotheses:
\begin{itemize}
  \item The modality-specific classifier approach (pre-filtering audio before it
        reaches the language backbone) is architecturally straightforward and
        should be the first defence evaluated, given its success in the image
        domain.
  \item Session-level context tracking (established as necessary by the Crescendo
        finding) applies directly to audio conversations; whether the same
        multi-turn priming signature is detectable in acoustic features is
        unknown.
\end{itemize}

\paragraph{Evaluation infrastructure (blocking for all above)}
No standardised benchmark for audio or video safety evaluation exists.
A minimum viable benchmark would specify: a fixed set of audio-encoded attack
prompts covering at least the perturbation, steganographic, and multi-turn
categories above; a judge model capable of evaluating audio-language model outputs;
and a harm taxonomy adapted for audio-specific content (e.g., voice-cloning fraud,
real-time deepfake facilitation).
Without such infrastructure, empirical progress cannot be compared across systems
and the audio/video gap cannot be closed in a reproducible way.
This is identified as open problem MM3 (Section~\ref{sec:open_problems}).

\subsection{Industry Deployments}

Microsoft Azure Content Safety~\cite{azure2024contentsafety} provides a commercial
API covering text and image moderation.
It operates as a post-generation filter and does not publish latency or accuracy
figures in reproducible form.
Its multimodal coverage (text + image) is broader than most academic proposals,
but independent evaluation is absent.

Meta's Llama Guard series~\cite{meta2024llamaguard2} (v1, v2, v3) represents the
most complete documented progression of a safety classifier through multiple
iterations, including taxonomy evolution.
Llama Guard's text-only classification architecture is the most robustly verified
finding in this review; its extension to image inputs in later versions has not been
evaluated independently.
